\section{Weighted consequences}\label{sec:weighted-consequences}

The following corollary applies Theorem~\ref{thm:main} through Seacrest's
established equivalences of the corresponding assertions over all nonempty
finite oriented graphs~\cite{Seacrest2015}.  The equivalences themselves are
not new here.

\begin{corollary}[Weighted second-neighborhood forms]\label{cor:weighted-forms}
Let \(D\) be a nonempty finite oriented graph, and write \(\operatorname{Arc}(D)\) for its
arc set.  For a vertex weighting \(\eta:V(D)\to[0,\infty)\), write
\(\eta(S)=\sum_{x\in S}\eta(x)\).  Then the following hold.
\begin{enumerate}[label=\textup{(\roman*)}]
\item For every vertex weighting \(\eta:V(D)\to[0,\infty)\), there is a
vertex \(v\in V(D)\) such that
\[
 \eta(N_1^+(v))\leq\eta(N_2^+(v)).
\]

\item There is a vertex weighting \(\eta:V(D)\to[0,\infty)\) with
\(\sum_{x\in V(D)}\eta(x)=1\) such that
\[
 \eta(N_1^+(v))\leq\eta(N_2^+(v))\qquad\text{for every }v\in V(D).
\]
The entries of this weighting need not all be positive.

\item For every arc weighting \(w:\operatorname{Arc}(D)\to[0,\infty)\), extend the notation
by setting \(w(vs)=0\) when \(vs\notin \operatorname{Arc}(D)\).  For \(v,s\in V(D)\), put
\[
 \beta_v^w(s)=\max\Bigl(\{0\}\cup
 \{w(us)-w(vs):u\in V(D),\ v\to u\to s\}\Bigr).
\]
Then there is a vertex \(v\in V(D)\) such that
\[
 \sum_{s\in V(D)}\beta_v^w(s)
 \geq \sum_{u\in N_1^+(v)}w(vu).
\]
The maximum is zero if there is no indicated two-step path.  Such paths
are determined by the arcs of \(D\), so an arc of weight zero is still
present and is not treated as a missing arc.
\end{enumerate}
\end{corollary}

\begin{proof}
Seacrest's Proposition~2 proves that his arc-weighted Conjecture~4 is
equivalent to the unweighted conjecture, so Theorem~\ref{thm:main} gives
\textup{(iii)}.

To obtain \textup{(i)}, set $w(us)=\eta(s)$ on every arc.  For each $v,s$,
\[
 \beta_v^w(s)=
 \begin{cases}
  \eta(s),&s\in N_2^+(v),\\
  0,&s\notin N_2^+(v).
 \end{cases}
\]
Indeed, when $v\to s$, every eligible difference $w(us)-w(vs)$ is zero.
At an exact second out-neighbor, $w(vs)=0$ and the maximum is $\eta(s)$;
no other endpoint contributes.  Thus the two sides of \textup{(iii)}
become $\eta(N_2^+(v))$ and $\eta(N_1^+(v))$.  This is the
specialization noted immediately after Seacrest's Proposition~2.

For \textup{(ii)}, let $M_D$ have entry $1$ at $(v,s)$ when
$s\in N_2^+(v)$, entry $-1$ when $s\in N_1^+(v)$, and zero otherwise.
Reversing all arcs transposes this matrix:
$M_{\overleftarrow D}=M_D^{\mathsf T}$.  Seacrest's Theorem~1 applies
Farkas' lemma to give either a nonzero $\eta\ge0$ with $M_D\eta\ge0$,
or a $\xi\ge0$ with $M_D^{\mathsf T}\xi<0$ at every coordinate.
The latter contradicts \textup{(i)} on $\overleftarrow D$.
Normalize the former weighting to total weight one, proving \textup{(ii)}.
\end{proof}

\begin{remark}[A direct way to understand the vertex-weighted form]
There is also a direct route from the unweighted theorem to \textup{(i)}.
For nonnegative integer weights, replace vertex $x$ by an independent set
of size $\eta(x)$, putting all arcs between blocks prescribed by $D$.
If the total weight is zero the assertion is immediate; otherwise the
new graph is nonempty and oriented. A good clone of $v$ has exactly
$\eta(N_1^+(v))$ first neighbors. Its exact second neighbors lie in the
blocks of $N_2^+(v)$, and so there are at most $\eta(N_2^+(v))$ of them.
Consequently its unweighted inequality gives the desired weighted one.
The containment can be strict when a two-step path in $D$ has a zero-weight
intermediate vertex; only an upper bound is needed. Clear denominators for
rational weights. Finally approximate any nonnegative real weighting by
nonnegative rational weightings. Since $D$ has finitely many vertices,
some witness vertex occurs along an infinite subsequence. Taking limits
in its finite sums proves the real-weight assertion.

The two vertex-weighted forms have different witnesses. In \textup{(i)}
the vertex may depend on an arbitrarily specified weighting. In
\textup{(ii)} the weighting is chosen so that all vertices work at once.
Zero entries can be necessary: in the graph with the sole arc $a\to b$,
requiring both vertices to work forces $\eta(b)=0$, whereas
$\eta(a)=1$, $\eta(b)=0$ is a valid normalized solution.
\end{remark}
